\documentclass[11pt,a4paper]{article}
\usepackage[T1]{fontenc}
\usepackage{lmodern,amsmath,amssymb,amsthm,mathtools,booktabs,array,longtable}
\usepackage[a4paper,margin=27mm]{geometry}
\usepackage{microtype,enumitem}
\usepackage[hidelinks,hypertexnames=false]{hyperref}
\usepackage{xurl}
\setlength{\emergencystretch}{3em}
\newtheorem{theorem}{Theorem}
\newtheorem{lemma}[theorem]{Lemma}
\newtheorem{proposition}[theorem]{Proposition}
\theoremstyle{definition}
\newtheorem{assumption}[theorem]{Assumption}
\newcommand{\E}{\mathbb E}
\newcommand{\Pp}{\mathbb P}
\newcommand{\CE}{C_{\mathrm E}}
\newcommand{\ind}{\mathbf 1}
\title{Five summands in the iid Berry--Esseen conjecture\\
\large A sign-count proof and the finite-index dependency chain}
\author{Kacper Rodziewicz \and Bartosz Ko{\l}odziejek}
\date{Working draft --- 20 October 2026}
\begin{document}
\maketitle
\begin{center}\small
A revised and corrected version is intended for submission to arXiv.
\end{center}
\begin{abstract}
Using Prawitz's smoothing inequality and Shevtsova's iid estimates, we prove
\(C_5\le C_{\mathrm E}=(3+\sqrt{10})/(6\sqrt{2\pi})\) from the published
four-summand bound. A fixed-index extremal argument gives an attained
closed-CDF witness and bounded support before higher moments are used.
Complete moment and sign-probability covers reduce the remaining gapped
regions to exact rational count inequalities. A direct conditional-binomial
and Cantelli comparison treats the final corner without a support-gap,
Fourier or predecessor assumption. Endpoint atoms, closed seams and both
infinite threshold rays are retained throughout. The fifth-index argument
uses no eventual-in-\(n\) theorem. Under separately stated analytic and
historical continuation hypotheses, a finite chain also gives
\(C_n<C_{\mathrm E}\) for every \(n\ge6\); these hypotheses are kept explicit
in the dependency diagram. ChatGPT and Codex assisted the development of
the proofs and computational certificates and the preparation of this
manuscript, using the GPT-5.6 Sol, GPT-6.1 Sol and GPT-6 Astra models.
\end{abstract}

\section{Statement and precise scope}
For a real law $P$ with $\E X=0$, $\E X^2=1$ and $\rho=\E|X|^3<\infty$, let
\[
 R_n(P)=\frac{\sqrt n}{\rho}\sup_{t\in\mathbb R}
 \left|\Pp(X_1+\cdots+X_n\le t)-\Phi(t/\sqrt n)\right|,
 \qquad C_n=\sup_P R_n(P),
\]
where the $X_i$ are independent with the \emph{same} law $P$. Write
\[
 \CE=\frac{3+\sqrt{10}}{6\sqrt{2\pi}}
       =0.40973218\ldots.
\]
We prove the unrestricted assertion $C_5\le\CE$ from the four-summand input.
We use the published bound $C_4\le\CE$ from
\cite[Theorem~26.1, p.~174]{RK}. That theorem states the stronger
$C_4\le0.4097$; the stronger value is not needed here.

\begin{theorem}\label{thm:composition}
Assume $C_4\le\CE$. Then $C_5\le\CE$.
\end{theorem}

This is a computer-assisted proof. Its continuous inequalities, finite
certificates, and complete domain coverage are separate proof obligations.
The final verification record identifies the exact successful
proof objects, their domains, and their input files. The continuous
support, Fourier, resource and count arguments are supplied in the companion
proof sources indexed in Section~\ref{sec:sources}. This note explains the
composition and proves the new gap-free part in detail; it does not replace
those companion proofs by numerical success messages. No eventual-in-$n$ theorem or monotonicity of $C_n$ is used.

Besides $C_4\le\CE$, the cited inputs are Prawitz smoothing in the
normalization of \cite[Lemma~4.1, p.~262]{S2012} and the iid bounds of
\cite[p.~125]{S2013}. The continuous reductions and exact arithmetic certificates are supplied
in the accompanying proof supplement \cite{C5ProofSupplement}.

\section{Why a contrary law is a legitimate object}
\begin{lemma}[Fixed-index attainment and the correct jump edge]\label{lem:attain}
For every fixed integer $n\ge1$, $C_n$ is attained. After reflection of the
whole marginal if necessary, there are an attaining standardized law $P$
and a finite raw threshold $H$ such that
\[
 C_n=\frac{\sqrt n}{\rho}\{\Pp(S_n\le H)-\Phi(H/\sqrt n)\}>0.
\]
\end{lemma}
\begin{proof}
The symmetric two-point law gives a positive lower bound for $C_n$, while
$R_n(P)\le\sqrt n/\rho$. Consequently a maximizing sequence has bounded
third absolute moments. It is tight. On a weakly convergent subsequence,
first and second moments pass to the limit, since their tail expectations
are at most $M/L^2$ and $M/L$, respectively. If $\rho_j\to r$, lower
semicontinuity gives $\rho(P)\le r$. No uniform integrability of third
moments has been assumed.

Because $n$ is fixed, the $n$-fold product laws, and hence their sums,
converge weakly. All normalized sums have variance one, so near-maximizing
thresholds stay bounded by Chebyshev's inequality and the Gaussian tails.
For a convergent sequence of thresholds $z_j\to z$ the exact bounds are
\[
 \limsup F_j(z_j)\le F(z),\qquad
 \liminf F_j(z_j)\ge F(z-).
\]
Both right and left jump errors are bounded by the Kolmogorov distance from
the continuous Gaussian CDF. Thus
$\limsup\Delta_n(P_j)\le\Delta_n(P)$. It follows that
\[
 C_n r\le\sqrt n\Delta_n(P)\le C_n\rho(P)\le C_n r.
\]
Equality proves attainment, including equality of the limiting third
moments.

The functions $F(z)-\Phi(z)$ and $\Phi(z)-F(z-)$ are upper
semicontinuous and vanish at infinity. Their positive maxima are attained.
Reflection converts the second branch exactly into the first, because
$F_{-S}(-H)=1-F_S(H-)$. The endpoint atom is retained.
\end{proof}

Assume now that $C_5>\CE$, and choose the witness from Lemma~\ref{lem:attain}.
For this same marginal its predecessor score satisfies $R_4(P)\le C_4$.
The contact allowance is therefore
\begin{equation}\label{eq:qzero}
 q_5=5\bigl(R_4(P)-C_5\bigr)_+=0.
\end{equation}
Equation~\eqref{eq:qzero} is the only use of the four-summand input.
It does not assert $C_5\le C_4$.

\subsection{Contact before any higher moment is used}
Put $z=H/\sqrt5$, $F=\Pp(S_5\le H)$, $a=(F-\Phi(z))/\rho>0$,
$\beta=\E[X|X|]$, and $h(x)=\Pp(S_4\le H-x)$.
Contaminate the common marginal by an atom at an arbitrary real $x$ and
then center and scale the entire mixture. The raw event $S_5\le H$ is
kept fixed. Its first derivative is $5(h(x)-F)$, so atoms cause no CDF
differentiation issue. The central third moment derivative is
$|x|^3-\rho-3x\beta$, with the variance normalization adding
$-3\rho(x^2-1)/2$. Global maximality yields
\begin{equation}\label{eq:contact}
 5h(x)\le a|x|^3+q_0+q_1x+q_2x^2,
\end{equation}
where
\[
 q_0=5F+\tfrac12(a\rho+z\phi(z)),\quad
 q_1=-\sqrt5\phi(z)-3a\beta,\quad
 q_2=-\tfrac12(3a\rho+z\phi(z)).
\]
The nonnegative slack has expectation zero. It is lower semicontinuous,
so equality holds at every support point. Since the cubic coefficient
$a$ is positive and $0\le5h\le5$, the support is bounded. Higher moments
are used only after this conclusion.

Quantitative support follows by inserting the all-threshold predecessor
bound into this contact equality and using \eqref{eq:qzero}. Its numerical
certificate covers both support rays and the whole threshold line,
including both infinite threshold tails. A finite mesh without its
curvature remainder would not be sufficient.

\section{The complete moment-domain cover}\label{sec:cover}
Let $\beta=\E[X|X|]$. An elementary moment restriction is
\begin{equation}\label{eq:beta}
 \rho\ge\frac{1+\beta^2}{\sqrt{1-\beta^2}}.
\end{equation}
Indeed, if $m=\E X_+=\E X_-$ and $a_\pm=(1\pm\beta)/2$, signwise
Cauchy inequalities give $m\le\sqrt{a_+a_-}$ and
$\rho\ge(a_+^2+a_-^2)/m$. Define
\[
 b(r)=\left(\frac{r\sqrt{r^2+8}-r^2-2}{2}\right)^{1/2}.
\]
Every law with $\rho\le r$ satisfies $|\beta|\le b(r)$.
The certificates use the rational upper enclosures
\[
 b_2=\frac{84729828287999}{250000000000000}\ge b(237/200),
 \qquad
 b_3=\frac{35139006743001}{10^{14}}\ge b(6/5).
\]
These enclosing inequalities are checked exactly.

The low-moment cover is as follows. Each gapped row has a complete
100-panel partition of $p=\Pp(X>0)$; the last column is the number of
panels requiring count forests. Every other panel is excluded by its
acquired INITIAL constraints or a separately checked moment contradiction.
The beta scopes are signed and are not identified by reflection.

\begin{center}\small
\begin{tabular}{@{}llllr@{}}\toprule
Region & $\rho$ interval & $\beta$ interval & Method & Forests\\\midrule
B0 & $[1,11/10]$ & full & resources/counts &24\\
BN & $[11/10,23/20]$ & full & resources/counts &32\\
B2 S0 & $[23/20,467/400]$ & $[-b_2,1/5]$ & resources/counts &32\\
B2 S1 & $[23/20,467/400]$ & $[1/5,b_2]$ & resources/counts &18\\
B2 S2 & $[467/400,237/200]$ & $[-b_2,1/5]$ & resources/counts &34\\
B2 S3 & $[467/400,237/200]$ & $[1/5,b_2]$ & resources/counts &20\\
B3 S0 & $[237/200,477/400]$ & $[-b_3,1/5]$ & resources/counts &35\\
B3 T0 & $[237/200,477/400]$ & $[1/5,3/10]$ & resources/counts &22\\
B3 T1 & $[237/200,477/400]$ & $[3/10,b_3]$ & resources/counts &9\\
B3 S2 & $[477/400,6/5]$ & $[-b_3,1/5]$ & resources/counts &36\\
B3 T2 & $[477/400,6/5]$ & $[1/5,3/10]$ & resources/counts &22\\
B3 lower & $[477/400,6/5]$ & $[3/10,17/50]$ & resources/counts &11\\
B3 corner & $[477/400,6/5]$ & $[17/50,b_3]$ & gap-free &---\\\bottomrule
\end{tabular}
\end{center}
All intervals are closed. In particular the new boundary $17/50$ belongs
to both adjacent pieces. There are 295 gapped count forests in the
complete cover. Each retains all six original nearest-count roots.
All 1,770 roots are present, with 148,586 nodes and 75,178 leaves.
The other 905 panel instances have acquired resource exclusions.
The final verification record joins every row to its accepted proof objects.

For $\rho\in[6/5,7361/4000]$, 41 closed Fourier strips give an upper bound
below $4097/10000$. Their six support regions are
\begin{center}
\begin{tabular}{@{}lc@{}}\toprule
$\rho$ interval & Support radius $R$\\\midrule
$[6/5,5/4]$ & $13/5$\\
$[5/4,7/5]$ & $117/40$\\
$[7/5,3/2]$ & $25/8$\\
$[3/2,8/5]$ & $133/40$\\
$[8/5,17/10]$ & $7/2$\\
$[17/10,7361/4000]$ & $15/4$\\\bottomrule
\end{tabular}
\end{center}
The support has been acquired before these supported Fourier bounds are
applied. The final finite-strip maximum is less than $0.409674$.
For the remaining tail, the published iid estimate gives
\[
 R_5(P)\le\frac{3031}{10000}\left(1+\frac{323}{500\rho}\right)
 \le\frac{819}{2000}=0.4095\quad(\rho\ge7361/4000).
\]
The exact endpoint substitution and monotonicity join the infinite tail
without a missing boundary.

\section{A direct sign-count inequality without a support gap}
This section contains the new step. It applies to every standardized common
iid marginal in its stated moment region, with no extremality, contact,
bounded-support or predecessor premise.

Put $p=\Pp(X>0)$, $q=1-p$, $m=\E X_+=\E X_->0$, $v=pq$ and $h=v/m$.
Zero atoms are assigned to the nonpositive class. Set
\[
 B=\ind_{\{X>0\}},\qquad Y=p+hX,\qquad W=Y-B.
\]
Then $\E[W\mid B]=0$. The two conditional variances are exactly
\begin{align}
 v_0&=p^2\left(\frac{q(1-\beta)}{2m^2}-1\right),\label{eq:v0}\\
 v_1&=q^2\left(\frac{p(1+\beta)}{2m^2}-1\right).\label{eq:v1}
\end{align}
For $J=B_1+\cdots+B_5$, conditional independence gives
\[
 \Pp(J=j)=b_j(p):=\binom5j p^jq^{5-j},\qquad
 V_j:=\operatorname{Var}\Bigl(\sum_iW_i\mid J=j\Bigr)=(5-j)v_0+jv_1.
\]
Conditioning on a particular sign vector gives independent conditional
variables. Every vector with the same count $j$ gives the same sum law,
which justifies conditioning only on the count.

\begin{lemma}[Gap-free count bound]\label{lem:cantelli}
Define
\[
 T(V,d)=\begin{cases}1,&d\le0,\\ V/(V+d^2),&d>0.\end{cases}
\]
For every real $z$,
\begin{equation}\label{eq:count}
 \Pp\Bigl(\sum_{i=1}^5Y_i\le z\Bigr)
 \le Q(z):=\min\left\{1,\sum_{j=0}^5 b_j(p)T(V_j,j-z)\right\}.
\end{equation}
Consequently, at $H=(z-5p)/h$, the normalized positive CDF error is at most
\begin{equation}\label{eq:score}
 \frac{\sqrt5}{\rho}
 \left[Q(z)-\Phi\left(\frac{(z-5p)m}{\sqrt5\,pq}\right)\right].
\end{equation}
\end{lemma}
\begin{proof}
A centered random variable $U$ with variance $V$ obeys
$\Pp(U\le-d)\le V/(V+d^2)$ for $d>0$. For $V>0$ this follows by
applying Markov's inequality to $(U-V/d)^2$ on that event; for $V=0$
the event is empty. At $d=0$ the safe bound is one, including $V=0$.
Apply this separately to each conditional sum and use its actual binomial
probability. Finally $\sum Y_i=5p+hS_5$ gives the normal coordinate.
\end{proof}

The convention $T(0,0)=1$ is essential. Substituting zero would omit a
possible atom exactly on an integer count boundary.

\subsection{Reducing the corner to two-dimensional boxes}
The additional moment constraints are
\begin{equation}\label{eq:moments}
 m\ge\frac{1+\beta^2}{2\rho},\qquad
 \frac{2m^2}{1+\beta}\le p\le1-\frac{2m^2}{1-\beta},\qquad m^2\le pq.
\end{equation}
The first two follow from the signwise Cauchy inequalities already used.
For the last, conditional second moments give
$1\ge m^2/p+m^2/q=m^2/(pq)$.

Consider a closed beta slab $[b_l,b_u]$ in $[17/50,b_3]$, keeping the
whole moment interval $\rho\in[477/400,6/5]$. Set
\[
 m_l=\frac{1+b_l^2}{2(6/5)},\quad
 p_l=\frac{2m_l^2}{1+b_u},\quad
 p_u=1-\frac{2m_l^2}{1-b_l}.
\]
Every law lies in this $p$ interval; all the selected intervals are contained
in $(0,1/2)$. For a closed subbox $p\in[l,u]$, $z\in[a,b]$, use
\begin{align*}
 \bar v_0&=\max\left\{0,u^2\left(\frac{(1-l)(1-b_l)}{2m_l^2}-1\right)\right\},\\
 \bar v_1&=\max\left\{0,(1-l)^2\left(\frac{u(1+b_u)}{2m_l^2}-1\right)\right\}.
\end{align*}
These bound \eqref{eq:v0}--\eqref{eq:v1} above: each actual bracket is
nonnegative, and each positive factor has been enlarged. A negative
upper bracket means that the actual-law domain is empty; clipping it to
zero remains safe. The exact upper binomial mass is attained at
$p_j=\min(u,\max(l,j/5))$. Hence an upper CDF on this whole box is
\[
 \bar Q=\min\left\{1,\sum_{j=0}^5 b_j(p_j)
 T\bigl((5-j)\bar v_0+j\bar v_1,j-b\bigr)\right\}.
\]

Let $\nu=a-5u$. A lower bound for the normal argument is
\begin{equation}\label{eq:normal}
 x_l=\begin{cases}
 \nu m_l/(\sqrt5\,u(1-u)),&\nu\ge0,\\
 \nu/\sqrt{5u(1-u)},&\nu<0.
 \end{cases}
\end{equation}
The first branch uses the lower mean and upper variance factor. In the
second branch use $m\le\sqrt{pq}$; if the actual numerator is positive
the negative lower bound is automatic. Otherwise note that for $z\ge0$
and $p<1/2$,
\[
 \frac{d}{dp}\frac{z-5p}{\sqrt{p(1-p)}}
 =-\frac{5p+z(1-2p)}{2[p(1-p)]^{3/2}}<0.
\]
Thus \eqref{eq:normal} holds on the entire box. When evaluating it exactly,
the negative numerator is divided by a \emph{lower} square root. Reversing
that rounding direction would invalidate the bound.

The resulting whole-box score bound is
\[
 \frac{\sqrt5}{477/400}\max\{0,\bar Q-\Phi(x_l)\}.
\]
Taking the positive part before multiplication avoids an incorrect
upper-factor substitution on a negative difference.

\subsection{The finite certificate and the infinite phase rays}
The twelve closed beta slabs have endpoints
\[
 0.340,0.341,\ldots,0.351,b_3.
\]
For each slab, the $p$ interval is rounded outward to a $10^{-6}$ grid.
The five phase intervals $[0,1],\ldots,[4,5]$ give 60 initial boxes.
Recursive midpoint splits in $p$ and $z$ produce a complete forest with
1,994 nodes and 1,027 leaves. Exact acceptance recomputes every leaf and
checks both children of every split. Its maximum stored score upper is
\[
 \frac{20484834187}{50000000000}=0.40969668374<0.4097<\CE.
\]
A separately implemented rational Gaussian evaluator checks every leaf;
its maximum is $204848338979/500000000000$, no larger than the stored
bounds. It does not import the original numerical oracle.

For completeness let $m_L,p_L,p_U$ denote the same moment bounds for the
\emph{whole} corner $\beta\in[17/50,b_3]$. Set
\[
 v_U=p_U(1-p_U),\quad s_U=v_U^2/m_L^2-v_U,\quad
 V^U_j=\max\left\{\frac j{p_L},\frac{5-j}{1-p_U}\right\}s_U.
\]
Here $s=\E W^2=v^2/m^2-v\le s_U$: on $v\ge m_L^2$ the function
$v^2/m_L^2-v$ is increasing. Also
$b_j(p)\le\binom5j p_U^j(1-p_L)^{5-j}$.
For $z\le0$, drop the nonnegative normal CDF and apply monotonicity at
zero. This gives the uniform bound
\[
 \frac{\sqrt5}{477/400}
 \left[(1-p_L)^5+\sum_{j=1}^5\binom5j p_U^j(1-p_L)^{5-j}
 \frac{V^U_j}{V^U_j+j^2}\right]<0.331.
\]
For $z\ge5$, $F\le1$ and the normal argument at five is
$\sqrt5\,m/p$. Cantelli for the standard normal gives the bound
\[
 \frac{\sqrt5}{477/400}\frac1{1+5m_L^2/p_U^2}<0.187.
\]
Both displayed inequalities are verified over exact rationals with an
outward square root of five. The two rays and the compact closed interval
cover every real threshold.

\begin{proposition}[Independently certified gap-free corner]
For every standardized iid marginal with
\[
 \rho\in[477/400,6/5],\qquad\beta\in[17/50,b_3],
\]
and every raw threshold $H$, the normalized positive CDF error is strictly
less than $4097/10000$. No $C_4$, support, contact, or physical-gap assumption
is needed for this proposition.
\end{proposition}

\section{What the other certificates prove}
The remaining low-moment regions concern the attained contrary witness.
Their chain is
\[
\begin{aligned}
 \text{contact and }q_5=0
 &\;\Longrightarrow\;\text{support and Fourier bounds}\\
 &\;\Longrightarrow\;\text{INITIAL}\;\Longrightarrow\;\text{resources}
 \;\Longrightarrow\;\text{counts}.
\end{aligned}
\]
Each arrow has a continuous proof as well as a source-specific arithmetic
obligation. INITIAL is a finite induction started with \emph{no} physical
defect or gap seed. All surviving final cells must acquire a genuine
central support gap; repeated improvement without that gate proves
nothing about a gap.

The defect used in this induction is
\[
 D=(3+\sqrt{10})\rho-\gamma-3J,\qquad
 \gamma=\E X^3,\quad J=\sum_{(a,b)}\tau_{a,b}(b-a)^2,
\]
where the sum runs over bounded gaps of the support and
$\tau(t)=\E[X\ind_{\{X\ge t\}}]$ is constant on each gap.
The independently derived conserved identity comes from integrating the
same contact slack against $\tau$ in two ways and using compensated smooth
Fourier inversion. It includes the complete high-frequency tail and
keeps the half-open gap convention at atoms. Its proof does not use an
asymptotic expansion in $n$.

Once a central gap $(-A,B)$ is acquired, the resource induction keeps
bounds on the same marginal's moments, endpoint locations, and conditional
noise. One small but necessary repair retains the already available price
\[
 6mAB\ge6m_L a b>0\qquad(m\ge m_L,\ A\ge a,\ B\ge b).
\]
It can be combined with another lower price by taking their maximum.
Since $4AB\le(A+B)^2$, a lower price $P$ gives
$A+B\ge\sqrt{2P/(3m_U)}$ when $m\le m_U$.
The new term alone satisfies
\[
 \frac{2(6m_Lab)}{3m_U}\le4ab\le(a+b)^2.
\]
It therefore retains a gap already proved; it does not manufacture a
physical gap from a failed INITIAL calculation.

\subsection{The actual iid count probability}
For the same representation as in Section~4 and a nearest-count phase
$z=k+u$, let $W_j$ denote the sum of the five noises conditional on
exactly $j$ positive signs. The closed CDF obeys
\begin{equation}\label{eq:outgoing}
 \Pp(S_5\le H)\le F_{\mathrm{Bin}(5,p)}(k)
 -b_k(p)\Pp(W_k>u)
 +\sum_{j>k}b_j(p)\Pp(W_j\le u-(j-k)).
\end{equation}
The outgoing event is \emph{strict}, the incoming event is closed. Atoms
on the event plane must not be subtracted twice.
Incoming bounds use Cantelli or exponential estimates. An outgoing
minorant follows from
\[
 G_a(x)=-x^4+3a^2x^2+2a^3x=x(x+a)^2(2a-x),\qquad a>0.
\]
It is nonpositive for $x\le0$ and has global maximum
$(9/4+3\sqrt3/2)a^4$. Therefore
\[
 \Pp(U>u)\ge\frac{\E G_a(U-u)}{(9/4+3\sqrt3/2)a^4}.
\]
Bounded support and common conditional moment constraints control the
expectation. The two conditional third moments remain common variables
through all constraints. Linear or positive-denominator affine-ratio
objectives are maximized over the resulting bounded polygon at its
vertices. This is a relaxation containing every actual law; feasibility
of a relaxed point is never asserted to produce a convolution law.

The exact forest tests include the coupled moment and physical-defect
constraints, phase polynomials and Gaussian enclosures. A leaf closes
only through a proved emptiness predicate or an upper score strictly below
the fixed rational lower enclosure of $\CE$. Every original nearest-count
root $k=0,\ldots,5$ is present. Its domain is derived from acquired mean,
scale and threshold bounds, so thresholds outside the original roots are
not silently discarded.

\section{Closing the contradiction}
\begin{proof}[Proof of Theorem~\ref{thm:composition}]
Suppose $C_5>\CE$. Lemma~\ref{lem:attain} supplies a positive closed-edge
maximizer, and \eqref{eq:qzero} supplies its zero predecessor allowance.
Its moments satisfy $\rho\ge1$ and \eqref{eq:beta}.

If $\rho\ge6/5$, the complete Fourier cover and the cited tail bound
exclude it. Otherwise it belongs to one of the closed rows in the
low-moment table. Every gapped row is excluded by the fully acquired
resource/count composition. The one remaining row is excluded directly
by the gap-free proposition. All seams are covered, including
$\rho=477/400$, $\rho=6/5$, and $\beta=17/50$. Thus no contrary
maximizer exists, and $C_5\le\CE$.
\end{proof}

The proof needs only a single fixed sample size. It avoids the long
local-domination and propagation chain used for eventual cutoffs. Combining
this result with any separate theorem for other indices requires checking
that theorem's own scope; it is not part of the present conclusion.


\section{The finite continuation to six}
\label{cut6:section}

Write \(G_n\) for the assertion \(C_n<\CE\), with the supremum in
\(C_n\) taken over all centered, unit-variance laws having a finite
third absolute moment. The finite continuation \cite{Cutoff6Continuation}
establishes the implications below. Its analytic hypotheses are retained
explicitly; they are not consequences of the finite arithmetic.

\begin{assumption}[Retained finite-index inputs]
\label{cut6:ass:inputs}
Let \(n\in\{6,7,8,9\}\).
\begin{enumerate}[label=(\roman*),leftmargin=*]
\item \(\mathcal A_6\): the inherited LOW, physical-gap, resource,
charged-score and count reductions hold with their stated same-law
hypotheses. In particular, their finite acquisitions apply to an
unrestricted maximizing witness on the contrary branch \(C_n\ge\CE\).
The initial acquisition assumes no defect seed: each stage covers its
whole closed moment domain and uses only previously established bounds.
The resulting probability panels and all \(n+1\) original count roots
cover every surviving witness, and the acquired count majorant bounds
that witness's normalized score.
\item \(\mathcal F\): the smoothing inequalities, profile bounds and
elementary enclosures used in the supported Fourier estimates hold
uniformly on their stated domains, with their stated remainder bounds.
This includes the literature inputs, not merely evaluations at selected
parameter values.
\end{enumerate}
The separate tail hypothesis \(\mathcal T_{10}\) is
\[
 C_n<\CE\qquad\text{for every integer }n\ge10.
\]
Here \(\mathcal T_{10}\) denotes the earlier conditional continuation,
including its retained analytic and historical certificate premises.
\end{assumption}

Fixed-index attainment, reflection at an atom, the contact inequality and
the continuous support criterion have separate analytic proofs. They
are distinguished from the retained LOW and count reductions in
Assumption~\ref{cut6:ass:inputs}. The coarse global estimate
\begin{equation}
 C_5<\frac{429}{1000}
 \label{cut6:eq:coarse5}
\end{equation}
is acquired independently of \(G_5\), within the same analytic and
Fourier framework. No sharp fifth-index conclusion is used to initiate
the continuation.

\begin{theorem}[Conditional finite continuation]
\label{cut6:thm:finite}
Under \(\mathcal A_6\) and \(\mathcal F\),
\[
 G_6\wedge G_7\wedge G_8\wedge G_9.
\]
Under the additional hypothesis \(\mathcal T_{10}\),
\[
 C_n<\CE\qquad(n\ge6).
\]
\end{theorem}

\begin{proof}
Suppose first that \(C_6\ge\CE\). Attainment and whole-law reflection
give a maximizing law and a positive CDF branch at a finite raw threshold.
Reflection interchanges the right CDF with the corresponding left limit;
atoms require no continuity assumption. All subsequent acquisitions
refer to this one unrestricted maximizing law.

The contact allowance is licensed by \eqref{cut6:eq:coarse5}:
\[
 6\bigl(R_5(P)-C_6\bigr)_+
 <6\left(\frac{429}{1000}-\frac{4097}{10000}\right)
 =\frac{579}{5000},
 \qquad \frac{4097}{10000}<\CE.
\]
Use \(q_6=579/5000\) in the support criterion. Its full-ray estimates
are uniform in the raw threshold. The nine acquired closed moment ranges
cover \([1,1.84025]\); support is established before each supported
Fourier estimate is applied.
The supported Fourier bounds exclude \(1.185\le\rho\le1.84025\)
by the ceiling \(4097/10000<\CE\). Shevtsova's inequality
\cite{S2013} gives, for \(\rho\ge1.84025\),
\[
 R_n(P)\le0.3031\left(1+\frac{0.646}{\rho}\right)
 \le0.4095<\CE.
\]
Thus the remaining closed bands, with support acquired before its use,
are
\begin{center}
\begin{tabular}{@{}ccc@{}}
\toprule
Band&Third moment&Support\\
\midrule
\(B_0\)&\([1,11/10]\)&\(|X|\le23/10\)\\
\(B_N\)&\([11/10,23/20]\)&\(|X|\le12/5\)\\
\(B_2\)&\([23/20,237/200]\)&\(|X|\le5/2\)\\
\bottomrule
\end{tabular}
\end{center}
Every moment seam belongs to both adjacent closed domains.

Within these bands the complete INITIAL acquisitions yield the physical
gap and resource bounds. Each of the 100 closed probability panels is
then excluded either by a whole-panel contradiction or by complete
count forests. Every original root is retained; subdivisions cover both
closed children. Each terminal score bound is strictly below a fixed
rational lower bound for \(\CE\). Assumption~\ref{cut6:ass:inputs}
therefore excludes the maximizing witness. Hence \(G_6\).

Proceed successively at \(n=7,8,9\). On the contrary branch
\(C_n\ge\CE\), the already proved global predecessor \(G_{n-1}\)
gives, for the same maximizing law,
\[
 n\bigl(R_{n-1}(P)-C_n\bigr)_+=0.
\]
Use \(q_n=0\) in that index's own support acquisition, Fourier bound,
INITIAL acquisition and count cover. The preceding argument applies
separately at each index. It includes equality \(C_n=\CE\), uses no
monotonicity in \(n\), and proves \(G_7,G_8,G_9\), in that order.
Attaching \(\mathcal T_{10}\) proves the second assertion.
\end{proof}

\subsection*{An elementary compression of the finite cover}

\begin{lemma}[Sign probability]
\label{cut6:lem:sign}
If \(\E X=0\), \(\E X^2=1\), and \(\rho=\E|X|^3<\infty\), then
\[
 \Pp(X>0)\rho^2,\ \Pp(X<0)\rho^2
 \ge\kappa:=\frac8{27}(7\sqrt7-17)>\frac{11261}{25000}.
\]
\end{lemma}
\begin{proof}
Put \(p=\Pp(X>0)\), \(m=\E X_+=\E X_->0\), and
\(a=\E X_+^2\in(0,1)\). Cauchy on each sign gives
\[
 m^2\le pa,\qquad a^2+(1-a)^2\le m\rho,
\]
and consequently
\[
 p\rho^2\ge\frac{(2a^2-2a+1)^2}{a}.
\]
The derivative of the expression on the right is
\((2a^2-2a+1)(6a^2-2a-1)/a^2\); its minimum is attained at
\(a=(1+\sqrt7)/6\) and equals \(\kappa\). Reflection gives the
other sign. An atom at zero does not affect either inequality.
\end{proof}

For \(\rho\le r\), this confines \(p=\Pp(X>0)\) to
\[
 \frac{11261}{25000r^2}\le p\le1-\frac{11261}{25000r^2}.
\]
Strict endpoint comparisons discard whole closed probability panels.
The residual resource exclusions have two terminal forms. In the first,
acquired nonnegative coefficients and a physical defect give
\[
 D\ge d\ge a(1-\theta)^{3/2}+b\theta,
 \qquad 0\le\theta\le1,\qquad a>D,\quad b\ge\tfrac32a.
\]
Convexity makes the right side at least
\(a+(b-3a/2)\theta\ge a>D\), a contradiction. In the second,
the acquired rational bounds \(\rho_-\le\rho\le\rho_+\) satisfy
\(\rho_->\rho_+\). The complete four-index cover is
\begin{center}
\begin{tabular}{@{}lrrrrr@{}}
\toprule
&Sign&Tangent&Inconsistent \(\rho\)&Count forests&Total\\
\midrule
\(n=6,7,8,9\)&824&52&40&284&1200\\
\bottomrule
\end{tabular}
\end{center}
The count forests have 2,362 original roots and 257,089 leaves.
This compression removes redundant exposition, not the acquisitions
used inside the surviving resource and count arguments.

\subsection*{Dependency boundary}
The sharp fifth-index argument and the finite continuation have different
roles. Their logical arrangement is
\[
\begin{gathered}
\text{published inequalities}+\mathcal A_6,\mathcal F
 \Longrightarrow C_5<429/1000\\
\Downarrow\\[-2pt]
G_6\Longrightarrow G_7\Longrightarrow G_8\Longrightarrow G_9\\[3pt]
(G_6\wedge\cdots\wedge G_9)+\mathcal T_{10}
 \Longrightarrow (G_n)_{n\ge6},\\[3pt]
\text{separately: fifth-index argument, with its own inputs}
 \Longrightarrow (C_5\le\CE).
\end{gathered}
\]
The right-hand chain uses only the coarse fifth-index seed, not the sharp
fifth-index bound.
Conversely, the fifth-index argument does not require this tail branch.
An eventual or finite-tail theorem cannot replace a missing fifth-index
proof. Theorem~\ref{thm:composition} proves \(C_5\le\CE\), not
the strict assertion \(G_5\).


\section{Continuous reductions and arithmetic certificates}\label{sec:sources}
The proof of Theorem~\ref{thm:composition} consists of the following
implications:
\[
\begin{gathered}
 C_4\le\CE,\quad C_5>\CE
 \ \Longrightarrow\ \text{attained positive closed witness},\\
 \text{contact with }q_5=0
 \ \Longrightarrow\ \text{support before higher moments},\\
 \text{support and the same-law Fourier bounds}
 \ \Longrightarrow\ \text{finite moment and probability covers},\\
 \text{gapped count bounds or the direct gap-free comparison}
 \ \Longrightarrow\ \text{contradiction}.
\end{gathered}
\]
The direct gap-free proposition is valid for arbitrary laws in its
moment rectangle and has no incoming contact or support edge.
The conditional tail in Section~\ref{cut6:section} has no incoming edge
into this fifth-index proof.

The companion proof supplement \cite{C5ProofSupplement} supplies the
continuous support criterion, compensated Fourier inversion and
conserved identity, the complete LOW and HIGH inequalities, the finite
resource induction and the conditional count majorants used in the
gapped rows. These derivations are part of the proof, not consequences
of a successful numerical calculation. Every quantitative acquisition
is tied to the same unrestricted contrary witness and to its original
closed parameter domains.

For each accepted count forest, exact acceptance checks every initial
root, both children of every split and every terminal inequality.
A terminal region is either necessarily empty or has a rational score
upper bound below a fixed rational lower enclosure of $\CE$.
Thus the finite trees prove inequalities over continuous domains, not
over sampled points. Square roots are enclosed by integer squaring.
Gaussian, exponential and trigonometric values use finite series with
explicit remainder bounds. The gap-free corner has, in addition, a
separate rational evaluation of every leaf and both unbounded rays.

The mathematical derivations, source code and complete finite
certificates are available in the
\href{https://github.com/rodziewiczk/berry-esseen-045-certificate/tree/main/moment-defects}
{accompanying repository}.
The two archived source packages \cite{C5ProofSupplement,Cutoff6Continuation}
retain their distinct theorem scopes. In particular the conditional
tail hypotheses of Assumption~\ref{cut6:ass:inputs} are not discharged
merely by adjoining the fifth-index certificate.

\section*{AI assistance declaration}
ChatGPT and Codex were used in developing the mathematical arguments
and proofs, constructing and checking computational certificates, and
drafting and revising the manuscript. This included the GPT-5.6 Sol,
GPT-6.1 Sol and GPT-6 Astra models. Kacper Rodziewicz and Bartosz
Ko{\l}odziejek remain responsible for the mathematical claims and final text.

\begin{thebibliography}{9}
\bibitem{RK} K.~Rodziewicz and B.~Ko{\l}odziejek,
\emph{Moment defects and sharp iid Berry--Esseen bounds},
Hexagon 2610.00147, version~3. Four-summand input:
Theorem~26.1, p.~174.
\url{https://hexagonmath.org/2610.00147v3}.
\bibitem{S2012} I.~Shevtsova,
\emph{Moment-type estimates with asymptotically optimal structure for the
accuracy of the normal approximation},
Annales Mathematicae et Informaticae \textbf{39} (2012), 241--307.
Prawitz smoothing: Lemma~4.1, p.~262.
\url{https://publikacio.uni-eszterhazy.hu/3231/1/AMI_39_from241to307.pdf}.
\bibitem{S2013} I.~G.~Shevtsova,
\emph{On the absolute constants in the Berry--Esseen inequality and its
structural and nonuniform improvements},
Informatika i ee Primeneniya \textbf{7}(1) (2013), 124--125.
\url{https://www.mathnet.ru/eng/ia252}.
\bibitem{C5ProofSupplement} K.~Rodziewicz and B.~Ko{\l}odziejek,
\emph{The fifth iid Berry--Esseen constant: continuous reductions and
exact finite certificates}, proof supplement, 20 October 2026.
\url{https://github.com/rodziewiczk/berry-esseen-045-certificate/releases/tag/cutoff5-sources-2026-10-20}.
\bibitem{Cutoff6Continuation} K.~Rodziewicz and B.~Ko{\l}odziejek,
\emph{Conditional finite continuation at indices six through nine},
proof supplement, 20 October 2026. The retained analytic and tail
hypotheses are stated in Section~\ref{cut6:section}.
\url{https://github.com/rodziewiczk/berry-esseen-045-certificate/releases/tag/cutoff5-sources-2026-10-20}.
\end{thebibliography}
\end{document}
